% Generado por sistema/backend ReporteTex.java -- no editar a mano.
% Datos: dataset de prueba (respuestas sinteticas)
Con $a = 1\,664\,525$, $c = 1\,013\,904\,223$, $m = 2^{32} = 4\,294\,967\,296$ y semilla $x_0 = 12\,345$, las primeras iteraciones son:

\begin{mitabla}{Primeras iteraciones del generador congruencial mixto}
\begin{tabular}{crrrc}
\toprule
$n$ & $x_n$ & $a \cdot x_n + c$ & $x_{n+1} = (a\,x_n + c) \bmod m$ & $R = x_{n+1}/m$ \\
\midrule
0 & 12\,345 & 21\,562\,465\,348 & 87\,628\,868 & 0.0204 \\
1 & 87\,628\,868 & 145\,861\,455\,411\,923 & 71\,072\,467 & 0.0165 \\
2 & 71\,072\,467 & 118\,302\,912\,037\,398 & 2\,332\,836\,374 & 0.5432 \\
3 & 2\,332\,836\,374 & 3\,883\,065\,479\,336\,573 & 2\,726\,892\,157 & 0.6349 \\
4 & 2\,726\,892\,157 & 4\,538\,981\,181\,534\,648 & 3\,908\,547\,000 & 0.9100 \\
5 & 3\,908\,547\,000 & 6\,505\,875\,209\,079\,223 & 483\,019\,191 & 0.1125 \\
\bottomrule
\end{tabular}
\end{mitabla}

Por ejemplo, para $n = 0$: $a \cdot x_0 + c = 1\,664\,525 \times 12\,345 + 1\,013\,904\,223 = 21\,562\,465\,348$; al dividir entre $m$ el residuo es $x_1 = 87\,628\,868$ y $R = 87\,628\,868 / 4\,294\,967\,296 = 0.0204$. Estos son los valores $R_1, R_2, \ldots$ que usa la simulación.

